\scanpage{156}\mspage{30}
\unsure{dont|alors} les \unsure{quatre} premières utilisent seulement l'hypothèse $\Phi\longleftrightarrow\Psi$ (de a) de la définition 4.1.), et la dernière seulement, $d\Longrightarrow c'$, fait appel à l'hypothèse de factorisation de Quillen (b) dans 4.1.); et encore il suffit de la supposer pour $j$ seulement !.

$c')\Longrightarrow c)$ tautologie pure, $c)\Longrightarrow a)$ car on sait que $\widetilde\Phi$ est stable par facteurs directs (et contient $\Phi$), $a\Longrightarrow b$ car $\widetilde\Phi\subset\Psi^*$ ($\widetilde\Phi$ est par définition la plus petite partie \unsure{close|clos} à g. de $\mathrm{Fl}\,M$ contenant $\Phi$, et $\Psi^{\unsure{*|\wedge}}$ est close à gauche), $b\Longrightarrow d$ en \unsure{écrivant|réécrivant} le diagramme donné 4.2.2. (traits pleins) comme
\diagram{
\begin{tikzcd}[ampersand replacement=\&,column sep=large,row sep=large]
A\arrow[r,dash,"{\unsure{u}}","{\unsure{\phantom{\longrightarrow}}}"']\arrow[d,"j"',"{\in\Psi^*}"]\&X\arrow[d,"{p\in\Psi}"]\\
B\arrow[r,"v"']\arrow[ur,dashed,"q"']\&Y
\end{tikzcd}
\quad(\unsure{\alpha|\omega})
}
en y faisant \deleted{$\unsure{B=B'|X=B'}$} $Y=B$, $v=\mathrm{id}_B$ \inserted{et en \unsure{rapportant|remplaçant} $X$ \deleted{\unsure{et|par}} par $B'$, et $u$ par $i$ :} l'hypothèse b) $j\in\Psi^*$ signifie que dans un diagramme \inserted{traits pleins} $(\alpha)$, il existe tjrs $q$ telle que $qj=u$, $pq=v$, et en \unsure{l'exprimant|l'appliquant} dans le cas particulier où $v=\mathrm{id}_B$, on trouve d).

Reste à prouver $d)\Longrightarrow c')$, en utilisant cette fois l'hypothèse de factorisation sur $j$, i.e. qu'on peut écrire $j$ sous la forme $pi$ :
\diagram{
\begin{tikzcd}[ampersand replacement=\&,column sep=large,row sep=large]
A\arrow[rr,dash,"{\unsure{j|\Delta}}","{\unsure{\phantom{\longrightarrow}}}"']\arrow[dr,"{\unsure{\hat i|i}}"']\&\&B\arrow[dl,dashed,bend left=40,"q"]\\
\&B'\arrow[ur,"p"']\&
\end{tikzcd},
}
\unsure{avec} $i\in\Phi$, $p\in\Psi$. En vertu de d), \inserted{comme $p\in\Psi$} il existe un $q:B\longrightarrow B'$ \unsure{satisfaisant}
% archive sheet number: 155
